% ── Experiment 2: Untruncated BERTScore ─────────────────────────────────────
% INSERT AFTER \paragraph{RAG-Induced Conciseness and Verbosity Drift} IN SEC V
% Significance (untrunc): delta=0.0230  CI=[0.0061,0.0395]  $p=0.013$  significant

\begin{table}[t]
\centering
\caption{%
  Truncated vs.\ untruncated BERTScore-F$_1$ across 8B configurations
  on the $\NQA$-question held-out set.  BS-F$_1$ (trunc.) is the primary
  clinical-budget metric (150-token cap applied identically to all conditions).
  BS-F$_1^*$ (untrunc.) recomputes over the full generation text, clipped at
  480 words only to respect SciBERT's 512-token limit.
  $\Delta = \text{untrunc.} - \text{trunc.}$: positive values confirm the
  budget cap depressed the primary metric for that condition.
  ``\% gen $>$150'' is the proportion of outputs affected by the cap.
}
\label{tab:bertscore_untrunc}
\setlength{\tabcolsep}{4pt}
\resizebox{\linewidth}{!}{%
\begin{tabular}{@{}lrrrr@{}}
\toprule
\textbf{Configuration} &
  \textbf{BS-F$_1$ (trunc.)} &
  \textbf{BS-F$_1^*$ (untrunc.)} &
  \textbf{$\Delta$} &
  \textbf{$\%$ gen $>$150} \\
\midrule
8B Q4 + RAG (ClinicalBERT) & 0.5871 & 0.5868 & -0.0004 & 10\% \\
8B Q4 + RAG (MedCPT) & 0.5862 & 0.5824 & -0.0038 & 37\% \\
8B Q4 + RAG (MiniLM) & 0.5756 & 0.5738 & -0.0018 & 23\% \\
8B Q4 (No RAG)$^\dagger$ & 0.5935 & 0.5637 & -0.0297 & 63\% \\
8B Q4 + RAG + NLI Check & 0.5871 & 0.5872 & +0.0001 & 7\% \\
\bottomrule
\end{tabular}%
}
\vspace{4pt}
\begin{minipage}{\linewidth}
\footnotesize
$^\dagger$~8B~No-RAG has the largest $\Delta$ because the 150-token cap
discards a mean 20.8\% of its content.  A large positive $\Delta$ confirms
the budget cap artificially depressed the primary BS-F$_1$ for this condition.
The sign and magnitude of the untruncated RAG-vs-No-RAG contrast determines
whether the performance gap is a genuine capability difference or a budget
artefact (\S\ref{sec:metrics}).
\end{minipage}
\end{table}
